\section{Experimental Setup}
\label{sec:experiments}

\subsection{Data and study design}
We study ten-class jet tagging using the offline constituent view of the
JetClass dataset~\citep{JetClass}, introduced with Particle Transformer
(ParT)~\citep{Qu2022ParticleTransformer}.
The classes are QCD jets and nine signal categories:
$H\to b\bar b$, $H\to c\bar c$, $H\to gg$, $H\to4q$, $H\to q\bar q\ell\nu$,
$Z\to q\bar q$, $W\to q\bar q'$, fully hadronic top decay ($t\to bqq'$),
and leptonic top decay ($t\to b\ell\nu$).
All models receive the same reconstructed offline information.

The training population contains $1{,}000{,}000$ jets, balanced across the ten
classes. A disjoint, likewise balanced validation population contains
$125{,}000$ jets and controls stopping and checkpoint selection.
These populations are sampled without replacement from a 10M-event portion
of the training release. Evaluation uses the official
\emph{20-million-jet JetClass test sample}: $2{,}000{,}000$ jets per class, including
$2{,}000{,}000$ QCD jets, disjoint from the training and validation samples.

We organize the study as an eight-configuration matrix: standard versus selected
pair features, shared versus layerwise attention bias, and absence versus presence
of EdgeValue messages.
In configuration labels, \texttt{Std} denotes standard pair inputs and
\texttt{Sel} denotes standard inputs plus PT, TRACK, and REGION.
The following \texttt{S} or \texttt{L} denotes shared or layerwise bias,
and the suffix \texttt{+EV} identifies ParT-EV variants. Thus,
\texttt{Std-S} is the shared-bias baseline without EdgeValue.
Every configuration is trained with three random seeds, matched across
configurations, with the same training-event order for a given seed and epoch.
All gains are measured relative to the \texttt{Std-S} baseline trained under
the same protocol.

\subsection{Inputs and training}
Each jet is represented by at most 128 constituents in their original order,
with padding masked. The common 17-dimensional particle input
contains log transverse momentum and energy, their log fractions relative to
the jet, radial and signed angular offsets, charge, five
particle-identification indicators, and transverse/longitudinal impact parameters
with their uncertainties. Jet-relative quantities are calculated from the
constituent four-vectors. The selected relations use these same constituents
and their track measurements, without additional detector information.
Normalization constants are fitted on training data and shared across
configurations. Apart from the three factors under study, the backbone
architecture in Section~\ref{sec:method} is held fixed.

All models are trained from scratch with unweighted ten-class cross-entropy and
AdamW~\citep{Loshchilov2019AdamW}, using $(\beta_1,\beta_2)=(0.9,0.999)$,
weight decay $10^{-4}$, and a gradient-norm clipping threshold of 1.
Gradients are accumulated over two microbatches of 64 jets, giving an effective
batch size of 128. Training uses BF16 mixed precision; test inference uses FP32.

The training budget is 40 epochs, with a minimum of 12 epochs and an early-stopping
patience of eight epochs. The learning rate increases linearly to $10^{-3}$ and
then follows a cosine decay toward $10^{-5}$. Warm-up spans the first 5\% of
the planned optimizer updates; the schedule is fixed to the full 40-epoch
budget regardless of early stopping.

Checkpoints are selected by validation accuracy, using lower cross-entropy
to distinguish epochs within 0.01 percentage points of the best accuracy.
The same selection criterion governs early stopping for every configuration.

The comparison matches data, backbone dimensions, loss, optimizer, and training
budget, but does not impose equal parameter counts or equal training time.
Shared-bias models use sequence trimming, whereas layerwise-bias models do not;
their pair-processing paths also yield different BatchNorm populations.
Shared-versus-layerwise differences therefore cannot be attributed solely to
bias sharing. We assess EdgeValue by comparing models with the same feature
set and bias scheme, which retain the same trimming and pair-processing choices.

\subsection{Classification and background-rejection metrics}
We summarize classification performance with multiclass accuracy and macro
one-versus-rest area under the ROC curve (macro AUROC)~\citep{Fawcett2006ROC}.
For each trained model, class AUCs rank events by the corresponding softmax
probability; macro AUROC is their unweighted mean over the ten classes.
We report means and sample standard deviations across the three independently
trained models, rather than combining their predictions into an ensemble.
Accuracy and macro-AUROC differences from baseline are expressed in percentage
points.

Background rejection is measured separately for each signal class $s$ against
QCD. If $\ell_c(x)$ denotes the classifier logit for class $c$, we use
\begin{equation}
 \begin{aligned}
 D_s(x)&=\ell_s(x)-\ell_{\QCD}(x),\\
 \widehat\epsilon_{\QCD,s}(t)
   &=\frac{1}{n_{\QCD}}\sum_{x\in\mathcal Q}
      \mathbf{1}\!\left[D_s(x)\geq t\right],\\
 R_s(\epsilon_s)&=\frac{1}{\widehat\epsilon_{\QCD,s}(t_s)},
 \qquad \widehat\epsilon_{\QCD,s}(t_s)>0.
 \end{aligned}
 \label{eq:rejection}
\end{equation}
Here $\mathcal Q$ is the set of test QCD jets, $n_{\QCD}=|\mathcal Q|$,
and $\mathbf{1}[\cdot]$ is the indicator function; $t_s$ is the threshold at
target signal efficiency $\epsilon_s$, defined below.
The main rejection table follows the JetClass benchmark operating points:
50\% signal efficiency for the seven hadronic signal classes, 99\% for
$H\to q\bar q\ell\nu$, and 99.5\% for $t\to b\ell\nu$
\citep{Qu2022ParticleTransformer}. Each row states its signal efficiency;
the class-specific choices are the same for every configuration and seed.
A separate table reports rejection at 30\% efficiency for the seven hadronic
classes. All operating points use the same checkpoints and test predictions.

The discriminant is monotonic in $p_s/(p_s+p_{\QCD})$. At target efficiency
$\epsilon_s$, $t_s$ is the
$\lceil\epsilon_s n_s\rceil$-th largest score among the $n_s$ true signal events.
All events with $D_s\geq t_s$ pass, including ties; the achieved signal efficiency
can therefore exceed the target because of integer rounding or tied scores.
Thresholds are determined separately for each signal and trained model from
the full test sample.

Tables give the arithmetic mean and sample standard deviation of rejection
across the three training seeds. These standard deviations describe training
variability, not confidence intervals.
